Production system prompts stack many requirements, yet developers have little quantitative guidance on how these requirements compete. We ask how compliance with verifiable system-prompt instructions depends on how many there are, which ones are combined, where they are placed, and whether they are repeated. We run eight open-weight chat models (0.5B--8B parameters) on 5,120 controlled system-prompt configurations each and score every response with deterministic checks. Instruction count is the dominant factor. Mean per-instruction compliance falls from 83\% with one instruction to 70\% with five and 59\% with ten, and the share of responses that satisfy every instruction falls to 37\% and 5\%. Even the best model satisfies all ten only 24\% of the time. Other effects are much smaller than we hypothesized: placing an instruction last rather than first adds 1.7 percentage points, and repeating it adds 3.3. The small cost of repetition to other instructions matches the cost of a neutral filler line, so it comes from prompt length. Interference is specific rather than categorical. Instructions that fix the whole response layout, such as ``exactly 3 bullets'', often displace closing phrases, keywords, and self-introductions. Larger models comply more at every load, but they do not lose compliance more slowly per added instruction in absolute or odds terms. Practitioners should keep system prompts short, verify outputs with code, and not rely on ordering or repetition.
